Write $\angle B=3\beta$ and $\angle C=3\gamma$, so $\beta+\gamma=40^\circ$. The specified ordering gives $\angle PBC=2\beta$, $\angle PCB=2\gamma$, and $\angle EBC=\beta$, $\angle ECB=\gamma$. In triangles $BCE$ and $BCP$, the sine rule yields $\dfrac{BE}{BP}=\dfrac{\sin\gamma\sin100^\circ}{\sin140^\circ\sin2\gamma}=\dfrac{\sin100^\circ}{2\sin140^\circ\cos\gamma}$. Since $\sin100^\circ=\sin80^\circ=2\sin40^\circ\cos40^\circ$ and $\sin140^\circ=\sin40^\circ$, this becomes $\dfrac{\cos40^\circ}{\cos\gamma}$. Let $\theta=\angle BPE$. The sine rule in $BPE$ gives $\dfrac{BE}{BP}=\dfrac{\sin\theta}{\sin(\theta+\beta)}$. The value $\theta=50^\circ$ satisfies this because $\sin50^\circ=\cos40^\circ$ and $\sin(50^\circ+\beta)=\cos(40^\circ-\beta)=\cos\gamma$. Cross-multiplying and expanding the sine and cosine sums reduces the equation to $\sin\beta\,\sin(\theta-50^\circ)=0$. Since $\beta>0$ and $0<\theta<180^\circ-\beta$, the unique angle is $\boxed{50^\circ}$.
\begin{asy}
texpreamble("\usepackage{amsmath,amssymb}");
size(240);
defaultpen(fontsize(9pt));
pair B=(0,0),C=(1,0),A=(1,1/sqrt(3)),P=(.8263518223,.3007674664),E=(.7660444431,.1350743037);draw(A--B--C--cycle);draw(B--P--C);draw(B--E--C);draw(P--E,blue);label("A",A,N);label("B",B,SW);label("C",C,SE);label("P",P,NW);label("E",E,S);
\end{asy}
